\chapter{Harmonic-map singularities and concentration}\label{sec:examples-harmonic-energy}
These examples make the scale-invariant energy in Section \ref{sec:harmonic-map-energy} concrete and connect it with weak convergence of Radon measures. The harmonic-map and stereographic framework is described in \cite{Hel01}, Chapters 2--4; the energy, weak formulation and concentration calculations are given here.

\begin{example}[The radial projection is a singular stationary map]\label{ex:radial-harmonic-hedgehog}
For $m\ge3$, the map $u(x)=x/|x|$ on $\mathbb R^m\setminus\{0\}$, with any value assigned at zero, belongs to $W^{1,2}_{\mathrm{loc}}(\mathbb R^m;S^{m-1})$ and is a stationary weak harmonic map. It has
\[
 |Du|^2=\frac{m-1}{|x|^2},\qquad
 \Theta_u(0,r)=\frac{(m-1)|S^{m-1}|}{2(m-2)}.
\]
It has no continuous representative at the origin. Stationarity here does not assert energy minimization or stability.
\end{example}
\begin{proof}
For $r=|x|>0$, $Du=r^{-1}(I-u\otimes u)$, so its squared norm is $(m-1)r^{-2}$, its radial derivative is zero, and $\Delta u=-(m-1)r^{-2}u$. The energy is locally integrable precisely when $m>2$. Integration by parts for first derivatives on a punctured ball has an inner boundary term bounded by $C\varepsilon^{m-1}$, which tends to zero; thus the displayed classical derivative is the weak derivative across zero.

For the weak equation, integrate against a compactly supported smooth vector test outside $B_\varepsilon$. The normal derivative of $u$ on the inner sphere is zero, so the Laplacian integration has no inner boundary term. The integrable quantities pass to the limit and give the weak equation. For domain stationarity, the stress tensor is divergence-free away from zero. On a sphere its radial flux is $S\nu=e\nu$, with $e=(m-1)/(2\varepsilon^2)$. In the boundary integral against $X$, the constant term $X(0)$ integrates to zero by symmetry. Since $X(x)-X(0)=O(\varepsilon)$, the remaining flux is $O(\varepsilon^{m-2})\to0$. This proves stationarity across zero. Polar integration gives $E(B_r)=(m-1)|S^{m-1}|r^{m-2}/(2(m-2))$, proving the density. Distinct radial limits rule out a continuous representative.
\end{proof}

\begin{example}[The planar vortex fails the finite-energy hypothesis]\label{ex:planar-vortex-infinite-energy}
The map $u(x)=x/|x|$ from a punctured planar disk to $S^1$ is smooth and harmonic there, but it is not in $W^{1,2}$ on any disk containing the origin. On $\{\varepsilon<|x|<R\}$ its Dirichlet energy is $\pi\log(R/\varepsilon)$.
\end{example}
\begin{proof}
The same derivative formula as in Example \ref{ex:radial-harmonic-hedgehog}, now with $m=2$, gives $|Du|^2=r^{-2}$. Hence $E=\frac12\int_\varepsilon^R2\pi r^{-1}dr=\pi\log(R/\varepsilon)$, which diverges as $\varepsilon\downarrow0$. Altering the map at the single origin cannot change this integral. Thus the punctured-domain harmonic equation does not put this map within the finite-energy weak theory at the origin.
\end{proof}

\begin{example}[A spherical bubble concentrates its energy at one point]\label{ex:stereographic-energy-concentration}
Let $u$ be the stereographic map in Example \ref{ex:stereographic-cmc-reconstruction}, and put $u_\lambda(x)=u(x/\lambda)$ for $\lambda>0$. Each is smooth and harmonic on $\mathbb R^2$, and
\[
 e_\lambda(x)=\tfrac12|Du_\lambda(x)|^2=\frac{4\lambda^2}{(\lambda^2+|x|^2)^2},\qquad
 \int_{\mathbb R^2}e_\lambda=4\pi.
\]
As $\lambda\downarrow0$, the maps converge strongly in $L^2_{\mathrm{loc}}$ and weakly in $W^{1,2}_{\mathrm{loc}}$ to the south pole $(0,0,-1)$, whereas their energy measures converge weakly to $4\pi\delta_0$. Their gradients do not converge strongly in $L^2$ on a neighborhood of zero.
\end{example}
\begin{proof}
Scaling preserves the two-dimensional harmonic-map equation. The derivative formula in Example \ref{ex:stereographic-cmc-reconstruction} gives the displayed density, and polar integration gives
\[
 \int_{B_R}e_\lambda=8\pi\lambda^2\int_0^R\frac{r}{(\lambda^2+r^2)^2}dr
 =4\pi\frac{R^2}{\lambda^2+R^2}.
\]
The limit as $R\to\infty$ is $4\pi$. At every nonzero point $u_\lambda$ converges to the south pole. Its norm is one, so dominated convergence gives strong local $L^2$ convergence. The total gradient norm is uniformly bounded. Any weakly convergent gradient subsequence has distributional limit zero, determined by the strong convergence of the maps; weak compactness and uniqueness of this limit give weak convergence for the entire family on each relatively compact domain.

For a continuous compactly supported test function $f$, change variables $x=\lambda y$ to obtain
\[
 \int f(x)e_\lambda(x)dx=\int_{\mathbb R^2}f(\lambda y)\frac4{(1+|y|^2)^2}dy\longrightarrow4\pi f(0)
\]
by dominated convergence. The positive mass remaining in every $B_R$ prevents strong gradient convergence to zero. Thus the weak limit of the energy measures differs from the energy measure of the limiting map.
\end{proof}

\begin{example}[Higher-degree spherical bubbles]\label{ex:higher-degree-spherical-bubble}
Identify $\mathbb R^2$ with $\mathbb C$. For every integer $k\ge1$, the smooth map $v_k(z)=u(z^k)$, where $u$ is the preceding stereographic map, is harmonic and has energy density
\[
 e_k(z)=\frac{4k^2|z|^{2k-2}}{(1+|z|^{2k})^2},\qquad
 E(v_k;\mathbb R^2)=4\pi k.
\]
The rescalings $v_k(z/\lambda)$ concentrate energy $4\pi k$ at zero. For $k>1$ the differential vanishes at zero, but the map is smooth there.
\end{example}
\begin{proof}
Away from zero, $z\mapsto z^k$ is a local conformal coordinate change, so harmonicity follows from conformal invariance. The rational expression defining $u(z^k)$ is smooth at zero; consequently the harmonic-map equation extends there by continuity. The derivative of $z^k$ has conformal scale $k|z|^{k-1}$, proving the density. Polar integration and $s=r^{2k}$ give
\[
 E=8\pi k^2\int_0^\infty\frac{r^{2k-1}}{(1+r^{2k})^2}dr
 =4\pi k\int_0^\infty\frac{ds}{(1+s)^2}=4\pi k.
\]
The measure-concentration argument is the same change of variables and dominated convergence as in Example \ref{ex:stereographic-energy-concentration}, using the integrable density $e_k$. The power map has zero derivative at zero when $k>1$, which proves the last assertion. These explicit examples do not claim a general energy-quantization theorem for arbitrary weakly convergent harmonic maps.
\end{proof}
